\documentclass[]{../../../../tex/classes/styledArticle}
\usepackage{../../../../tex/packages/hebrewSupport}

\author{שחר פרץ}
\title{\textit{היסטוריה 40}}
\begin{document}
	\maketitle
	
	\subsection{בונים מדינה – ציר זמן}
	כל דבר שנאמר בכיתה – כדאי ללכת ולחפש איפה זה נמצא בספר. ככל שתארגנו את קטעי המקור בספר, יהיה לכם יותר קל בבחינה. 
	
	\begin{itemize}
		\item 1945 -- מלחמת העולם השנייה מסתיימת. לא רק היישוב היהודי בא''י בסערה. עוד מלפני שהמלחמה נגמרה, נפתחה המלחמה הקרה. חלוקת העולם אחרי המלחמה התרחשה בוועידה של סטלין, צ'רצ'יל ורוזוולט. אירופה הרוסה. נבחין בשתי בעיות עיקריות: 
		\begin{itemize}
			\item \textbf{בעיית העקורים: }אנשים שנאלצים להשאר במחנות, כי אין להם לאן ללכת, או שלא יכולים לחזור לבתיהם. רק בפולין, 1600 יהודים נרצחים בפוגרומים. בארה''ב הופעל לחץ כבד על הממשל האמריקאי לפתור את בעיית העקורים, ומכאן שמופעל גם לחץ על בריטניה ששולטת במקום. 
			
			\textbf{בעיית היהודים: }צ'רצ'יל לא נבחר בבחירות. העם רצה משהו חדש. מפלגת הלייבור הסוציאליסטית מתנערת מהבעיה. בריטניה חרבה וצריכה את הכסף האמריקאי. טרומן לוחץ על הבריטים, וועדה בסמכות טרומן מחליטה שצריך למצוא פתרון ל־100K יהודים באופן מיידי. וועידה משותפת, הוועדה האנגלו־אמריקאית, מגיעה לאותה המסקנה. למרות הוועדות הרבות, המדיניות הבריטית נותרה על קנה, ושיאה בנאום בווין בפרלמנט. מדיניות הספר הלבן השלישי (מ־1939) נמשכת. 
		\end{itemize}
		\item 1946 -- ההגנה בראשות בן־גוריון מבין שאנחנו בבעיה. כ־70\% מהיישוב היהודי תמך בהגנה, כ־20\% באצ''ל, וכ־10\% בלח''י. לאצ''ל וללח''י קראו ``הפרושים''. כאן, נכנסו לתקופה הקרויה ''תקופת המרי העברי``. לא מאהבת מרדכי אלא משנאת המן, התלכדו יחדיו המחתרות תוך הבנה שבריטניה לא תשנה את מדיניותה. 
		
		תוך כדי כך, אנשי הבריגדיה היהודית ויהודים אחרים באירופה, מנסים לארגן ספינות להעביר אנשים לבריטניה. אלו כל הזמן שולחים מבזקים לגבי המצב הרע באירופה. המחתרות האלו פעלו באופן לא חוקי. 
		
		על ועדת X עוד דיברנו בשיעור שעבר. השורה התחתונה היא שכל מחתרת פועלת בנפרד, אך עם תיאום. עם מטרה ברורה: להעיף את הבריטים מישראל, ולהקים מדינה יהודית. החשיבה העקרית היא פגיעה בתשתיות הבריטיות, כאמצעי לחץ. ואכן המשאבים להחזיק את א''י היו רבים. זהו \textit{מאבק צבאי}. שימו לב – יהיה עוד מאבק צבאי, שהתרחש לאחר פירוק תנועת המרי העברית. כל פעולות תנועת המרי הן פעולות צבאיות. בין הפעולות הבולטות נמנים ליל הגשרים, ליל הרכבות, פיצוץ תחנות רדאר ומשטרה, ליל המטוסים, והפריצה לעתלית. 
		
		סיבות לפירוק תנועת המרי: 
		\begin{itemize}
			\item כל הדברים לעיל הובילו ל''מבצע ברוסיד``. במסגרתו התרחש האירוע הראשון שהוביל לפירוק תנועת המרי – ''השבת השחורה``. שלושת אלפים (מתוך כמה מאות אלפים בכל היישוב היהודי) נעצרו, המוני מסמכים נאספו, וכמויות נשק הוחרמו. כבר בנקודת הזמן הזו, מתגבשת ההבנה בהגנה שפעילות צבאית נגד בריטניה לא משיגה את מטרותיה. 
			\item פיצוץ מלון המלך דוד: במשך שנים האצ''ל טען בתוקף שהיה לו אישור מועדת X. עוד לטענת האצ''ל הם התריאו. רק אחרי מלחמת ששת הימים תוקן החלק הדרום מערבי של המלון. 
			
			השורה התחתונה, היא שבעולם הפיצוץ מתפרש כפיגוע טרור. היישוב היהודי ומוסדות ההסתדרות הציונית מתנערים מהאירוע הזה. לימים, הסתבר שאכן קיבלו אישור מועדת X (\textit{הערה שלי: זה אי־דיוק. האישור התקבל לפיצוץ בתנאים אחרים}). 
		\end{itemize}
		
		בבגרות עדיף לכתוב קודם על השבת השחורה ואז על מלון המלך דוד. 
		\item 1946-כ''ט בנובמבר 1947 --
		כאן נכנסים לתקופה שהצורך של הפורשים להמשיך במאבק צבאי שלא בוחלים מפגיעה בחיי אדם, תנועת המרי מתפרקת, וכל מחתרת \textit{חוזרת} לדפוס המאבק הקודם שלה בבריטי. את הדפוסים הקודמים, אנו מחלקים לשניים: 
		\begin{itemize}
			\item \textbf{המאבק לעליה והתיישבות} (מאבק צמוד)\textbf{: }בהובלת ההגנה והמוסרות הציוניים בארץ ובעולם. מטרתו: פתיחת שערי א''י לעליה והתיישבות, סילוק המנדט הבריטי, והקמת המדינה. איך: אנחנו רואים באהדה הבינ''ל למצוקת העקורים מכנה רחב שיש עליו קונצנזוס עולמי נגד המדיניות הבריטית. נשתמש בזה כמנוף להשיג בדיפלומטיה, תוך כדי זה שנמשיך לעשות את כל המאמצים לקידום עליה בלתי־לגאליתהוהתיישבות (חומה ומגדל). 
			
			''אין לכם מושג כמה הבריטים משקיעים בלתפוס אוניות רעועות של מעפילים, רובן נשלחות לקפריסין. אבל יש לזה אימפקט בין־לאומי מאוד מאוד חשוב``. 
			\item \textbf{מאבק צבאי} (מאבק רצוף)\textbf{: }בהובלת האצ''ל והלח''י. מטרתו: סילוק המנדט הבריטי, פתיחת שערי א''י לעלייה יהודית, והקמת מדינה. איך: מטרת צבאי עד אחרון הבריטים בא''י. ידינו לא בוחלת בשום אמצעי – נפגע בחיילים, חפים מפשע, בריטים, ערבים, יהודים, נשדוד בנקים (גם בנקים יהודים) כדי לממן את הפעילות, העיקר נשיג את המטרה. 
		\end{itemize}
		אלו ואלו מנסים להשיג את אותן המטרות. אך התפיסה האידיאולוגית של הפורשים, אומרת שהבריטים עומדים בנינו לבין הקמת מדינה, והם האויב, ועל כן יש להלחם בהם. 
		
		מאוד אוהבים לשאול בבגרות מי עזר יותר להקמת המדינה. עצת המורה (לא שרית) היא לא לנקוט עמדה נרחצת, ולתת דוגמה לכל אחד מהמאבקים. זה חוקי. אם יש לכם שתי דוגמאות לאחד מהם, זה גם הולך. 
	\end{itemize}
	
	
	
	
	\ndoc
\end{document}